\section{Frobenius norms on adjoint groups of type \texorpdfstring{\(A_1\)}{A1}}\label{sec:adjoint}

Let \(G\) be the Mumford--Tate group of a polarized abelian variety
\(A/K\). Suppose every geometric simple factor of \(G^{\mathrm{ad}}\)
has type \(A_1\). The rational decomposition is
\[
 G^{\mathrm{ad}}=\prod_i\operatorname{Res}_{R_i/\Q}H_i,
\]
where \(H_i\) is an adjoint \(R_i\)-form of \(\mathrm{PGL}_2\).
Call an embedding \(R_i\hookrightarrow\overline{\Q}\) \emph{active} if
the Hodge cocharacter has nontrivial projection on the corresponding
geometric simple factor. Write \(b_i\) for the number of active
embeddings. On the adjoint Lie algebra a nontrivial projection has
weights \(-1,0,1\), because this Lie algebra is a subrepresentation of
\(\operatorname{End}(H^1)\). Each \(b_i\) is positive by the minimality
of the Mumford--Tate group.

\begin{theorem}[Adjoint trace bound]\label{thm:adjoint}
After a finite extension of \(K\), at a set of good absolute degree-one
primes \(v\) of natural density one, the following holds for every \(i\).
Let \(p\) be the residue characteristic. Above each prime
\(\mathfrak p\mid p\) of \(R_i\), all embeddings are active together
or inactive together, and their adjoint Newton slopes agree. Let
\(\mathcal F_i(v)\) consist of the active primes at which these slopes
are all zero. Then
\[
 \sum_{\mathfrak p\in\mathcal F_i(v)}
 f_{\mathfrak p}\bigl(1+\mathbf 1_{f_{\mathfrak p}\,\mathrm{odd}}\bigr)
 \le b_i-1,
 \qquad f_{\mathfrak p}=[(R_i)_{\mathfrak p}:\Q_p].
\]
At primes split completely in \(R_i\) there are at most
\(\lfloor(b_i-1)/2\rfloor\) active embeddings with flat adjoint Newton
slopes. If \(b_i\le 2\) for every \(i\), the ordinary primes have
natural density one over this extension.
\end{theorem}

For example, the two active embeddings of a rational adjoint factor
with a real quadratic label field can occur at separate split primes
or at one prime of residue degree two. In the first case determinant
parity makes a flat factor contribute two to the bound. In the second
case residue degree alone gives the same contribution. The argument
also applies to tensor representations, where one may have no
two-dimensional Frobenius block in \(H^1\).

\subsection{A common invariant of crystalline and \texorpdfstring{\(\ell\)}{l}-adic Frobenius}

Enlarge the base to define the absolute Hodge tensors of \(G\) and a
level for the corresponding Hodge-type Shimura datum, and to contain
its reflex field. Work at good
absolute degree-one primes with \(K_v=\Q_p\), away from \(2\), the level,
the discriminants of the \(R_i\), and the finitely many primes where
the datum lacks a hyperspecial model. At such a prime the abelian
scheme and its prime-to-\(p\) level give a point of the Siegel model.
The generic point extends to the Hodge-type model by finiteness of
the normalization of the closure. The crystalline tensor comparison
identifies the linear Frobenius and the Hodge filtration on a common
tensor fiber over \(\Q_p\). After trivializing over the completed
maximal unramified extension, write Frobenius as \(b\sigma\). Its class
satisfies the group Newton inequality and the Kottwitz equality for
the Hodge cocharacter by \cite[\S1.3 and Lemma~1.3.8]{kmps}.

At residue degree one, Kisin and Zhou's compatibility theorem
compares the class of linear crystalline Frobenius with the auxiliary
\(\ell\)-adic Frobenius classes, via isomorphisms respecting the
tensors over algebraic closures. More precisely, at good reduction
the compatible Weil--Deligne classes are unramified. Evaluation at
Frobenius gives a single Galois-stable semisimple \(G\)-class
\cite[\S5.3.3--5.3.4 and Theorem~5.3.10]{kzsemistable}. We use the
Frobenius convention giving slopes in \([0,1]\) on \(H^1\). The
invariant below is unchanged by inverting Frobenius.

On \(H_i\) take the regular function
\(T_i=1+\Tr(\operatorname{Ad}_{H_i})\), with the trace taken on its Lie
algebra. Evaluation on the Frobenius class yields \(T_i(v)\in R_i\):
the evaluations at the geometric factors form a Galois-equivariant
tuple indexed by embeddings of \(R_i\). If \(u,u^{-1}\) are the two
root eigenvalues, the invariant is \(2+u+u^{-1}\). They are ratios of
eigenvalues on \(H^1\), since the derived Lie algebra occurs inside
\(\operatorname{End}(H^1)\). The Weil bounds therefore give
\(0\le \tau(T_i(v))\le 4\) at all complex embeddings \(\tau\), and
integrality at all finite places away from \(p\).

\subsection{Determinant parity and the local valuation}

The centroid of the derived Lie algebra is \(\prod_i R_i\). Its
scalar action is an endomorphism of the Lie algebra as a Hodge tensor
representation, and descends with the other tensors. On the
crystalline fiber it commutes with linear Frobenius and preserves
the filtration. Thus this fiber of the adjoint Lie algebra decomposes
over the local fields \(L_{\mathfrak p}=(R_i)_{\mathfrak p}\).
On the rank-three \(L_{\mathfrak p}\)-summand the positive Hodge
filtration is either zero or a line. Its rank is the same at every
embedding of \(L_{\mathfrak p}\). By filtered comparison these are
precisely the inactive and active cases. The counts of active
embeddings are the Betti counts \(b_i\), since comparison identifies
the cocharacter classes and their action on the rational adjoint
quotients. The characteristic polynomial of linear Frobenius on this
summand is defined over \(L_{\mathfrak p}\), and gives the same slopes
at all its embeddings. The group Newton inequality bounds the
nonnegative adjoint slope by one at an active prime and by zero at an
inactive prime. In particular \(v_{\mathfrak p}(T_i(v))\ge0\) in the
inactive case, where the valuation is normalized by \(v_{\mathfrak p}(p)=1\).

In the active case the positive Hodge line becomes a nilpotent root
line over an algebraic closure. An inner form of \(\mathfrak{sl}_2\)
with a nonzero rational nilpotent is split: the trace-zero Lie
algebra of a division quaternion algebra has no such element. The
linear crystalline Frobenius acts on the split Lie algebra by an
automorphism, which is inner in type \(A_1\). Represent it by the
projective class of a lift \(B\in\mathrm{GL}_2(L_{\mathfrak p})\).
Over the completed maximal unramified extension, the summands are
indexed by embeddings \(\rho\) of \(L_{\mathfrak p}\). The coefficient
Frobenius \(\sigma\) permutes these embeddings cyclically, and the
linear part in the \(\rho\)-coordinate is \(\operatorname{Ad}(\rho(B))\).
The reference adjoint group is also split over \(L_{\mathfrak p}\),
because it is unramified of type \(A_1\). The transition to a tensor
trivialization preserves the centroid labels and the Lie bracket.
It therefore changes the tuple of projective matrices by
\(\sigma\)-conjugacy in the product of the split \(\mathrm{PGL}_2\)'s.
This tuple represents the projection of \([b]\) to the indicated
restriction of scalars.

The geometric fundamental group of \(\mathrm{PGL}_2\) is \(\Z/2\).
For a split factor the Kottwitz map is the lift's determinant
valuation modulo two, as follows by projecting the map for
\(\mathrm{GL}_2\). For the restriction of scalars the coinvariant
map sums these parities across the \(f_{\mathfrak p}\) embeddings.
Each matrix obtained from \(B\) has the same determinant valuation,
and each Hodge projection is the nontrivial minuscule class. The
Kottwitz equality therefore yields
\[
 f_{\mathfrak p}v_{\mathfrak p}(\det B)\equiv f_{\mathfrak p}\pmod 2.
\]

The nonnegative adjoint slope is the difference of the two eigenvalue
valuations of \(B\). If the difference is positive, each valuation
is integral by the quadratic Newton polygon over the unramified
field \(L_{\mathfrak p}\), so the difference must be one. The term
with smaller valuation determines the valuation of the trace, giving
\(v_{\mathfrak p}((\Tr B)^2/\det B)=-1\). If the eigenvalue valuations
coincide, they equal \(v_{\mathfrak p}(\det B)/2\). Their sum lies in
\(L_{\mathfrak p}\) and so has valuation at least the ceiling of this
number. Using the determinant parity when \(f_{\mathfrak p}\) is odd,
we obtain
\[
 v_{\mathfrak p}(T_i(v))
   =v_{\mathfrak p}\bigl((\Tr B)^2/\det B\bigr)
   \ge\mathbf 1_{f_{\mathfrak p}\,\mathrm{odd}}
   \quad(\mathfrak p\in\mathcal F_i(v)).
\]
The computation uses the linear Frobenius over \(\Q_p\), so the gap
one and the valuation \(-1\) hold without dividing a slope of a power
of Frobenius by \(f_{\mathfrak p}\).

\subsection{Bounded integer norms and ordinarity}

The local valuations bound the norm from below by
\[
 v_p(\Nm_{R_i/\Q}T_i(v))\ge -b_i+
 \sum_{\mathfrak p\in\mathcal F_i(v)}
 f_{\mathfrak p}\bigl(1+\mathbf 1_{f_{\mathfrak p}\,\mathrm{odd}}\bigr).
\]
If the sum reaches \(b_i\), the norm is integral at \(p\) and hence is
an integer in \([0,4^{[R_i:\Q]}]\). The same interval covers the case
of a zero norm. On \(G\) this norm is the regular class function
\(\prod_\tau T_{i,\tau}\). It is nonconstant: restrict to a torus in
one adjoint factor and take the identity in the others, giving
\(4^{[R_i:\Q]-1}(2+z+z^{-1})\).

Vasiu's theorem gives the Mumford--Tate conjecture here
\cite[Theorem~1.3.4(a) and Example~1.3.5(iv)]{vasiu}. Because the
auxiliary Galois action preserves the defining tensors, its Zariski
closure is the connected group \(G_{\Q_\ell}\). A proper conjugation-
invariant closed subset of this group cuts out a set of primes of
natural density zero by algebraic Chebotarev. Indeed the intersection
with the compact \(\ell\)-adic image has Haar measure zero, by Zariski
density of every open subgroup and the measure-zero property of the
zeros of a nonzero analytic function on each chart. Its shrinking
congruence neighborhoods give finite-quotient Chebotarev upper
bounds tending to zero. This is the bounded-trace density argument
used by Pink \cite[Lemmas~7.21--7.22]{pink}. Applying it to the
finitely many integer norm levels proves the inequality simultaneously.

Every flat active prime contributes at least two. Thus for all
\(b_i\le2\) the Newton and Hodge cocharacters have the same dominant
projection on every adjoint factor. On the torus \(G/G^{\mathrm{der}}\)
the Newton cocharacter is the Galois average of the Hodge cocharacter
by the Kottwitz equality. The Hodge conjugacy class is locally Galois
stable, since the degree-one completion contains the local reflex
field. At the unramified quasi-split group the dominant Hodge
representative is therefore Galois invariant. The adjoint and torus
projections jointly determine a rational cocharacter in the chosen
maximal torus, so the dominant Newton and Hodge cocharacters agree.
On \(H^1\) this means that all Newton slopes are zero or one.

\subsection{Exact half-slope multiplicities and a \texorpdfstring{\(p\)}{p}-rank bound}

For any geometric adjoint embedding in the rational factor \(i\),
let \(2c_i\) be the total dimension of the irreducible summands of
\(H^1_{\overline\Q}\) on which its simple Lie algebra acts nontrivially.
Dimensions are counted with multiplicity. This number is independent
of the embedding, since the representation is rational and Galois
acts transitively on the geometric factors above a rational factor.
Write \(r_i(v)=\sum_{\mathfrak p\in\mathcal F_i(v)}f_{\mathfrak p}\).

\begin{proposition}[Newton slopes and \(p\)-rank in adjoint type \(A_1\)]
\label{prop:adjoint-prank}
With the base extension of Theorem~\ref{thm:adjoint}, at every good
absolute degree-one prime outside the finite set used for the local
calculation, the Newton slopes are \(0,1/2,1\), with
\[
 \operatorname{mult}_v(1/2)=2\sum_i c_i r_i(v),\qquad
 f_p(A_v)=g-\sum_i c_i r_i(v).
\]
In particular, after also adjoining normal closures of the adjoint
label fields, on a set of primes of natural density one,
\[
 f_p(A_v)\ \ge\ g-\sum_i c_i\left\lfloor\frac{b_i-1}{2}\right\rfloor
 \ >\ \frac g2 \qquad (g>0).
\]
\end{proposition}

\begin{proof}
Over an algebraic closure an irreducible \(G\)-summand of \(H^1\)
is a character of the connected center times a tensor of symmetric
powers of standard \(\mathrm{SL}_2\)-modules, using the central
isogeny from the simply connected derived group times the center.
On an active factor the standard weights for the rational lift of
the Hodge cocharacter are \(-1/2,1/2\). The range between the largest
and smallest Hodge weights of the tensor is therefore the sum of
the symmetric-power exponents at active embeddings. Since \(H^1\)
has weights zero and one, this sum is at most one. A summand thus
either meets no active label or uses the standard module at exactly
one. In the second case its dimension is even. In particular
\(c_i\) is an integer and \(\sum_i c_i b_i\le g\), since distinct
active labels meet disjoint sets of irreducible summands.

The Newton and Hodge cocharacters have the same central component
under the rational lift, by the abelianization comparison above.
On every inactive factor both projections vanish. On each active
factor the Newton projection is either the Hodge projection or zero,
by the local calculation. Thus a summand using a flat active label
has the single Newton slope \(1/2\): its central weight is the
midpoint of its Hodge weights zero and one. Every other summand has
its original Hodge weights as Newton slopes. The \(r_i(v)\) flat
active labels in factor \(i\) account for \(2c_i r_i(v)\) copies
of \(1/2\). Polarization makes the multiplicities of zero and one
equal, giving the formula for \(f_p\).

At the completely split coefficient primes of the enlarged base,
Theorem~\ref{thm:adjoint} bounds \(r_i(v)\) by
\(\lfloor(b_i-1)/2\rfloor\). Since this integer is strictly smaller
than \(b_i/2\) for each \(i\) and \(\sum_i c_i b_i\le g\), the
bound is strict for positive-dimensional \(A\), including the case
of a torus where the sum is empty.
\end{proof}

Pink's ordinary-reduction theorem \cite[Theorem~7.1]{pink} assumes
\(\operatorname{End}(A_{\overline K})=\Z\) and treats a wider range
of root systems. Theorem~\ref{thm:adjoint} isolates the contribution
of each rational type \(A_1\) quotient and requires no endomorphism
hypothesis. For three active embeddings of one rational quotient the
inequality still allows one flat embedding at completely split
primes. In particular it leaves that case of cubic real multiplication
open. For comparison, Fit\'e's \(p\)-rank bound \(\lceil g/2\rceil\)
assumes an embedded Galois field of degree \(g\) in the geometric
endomorphism algebra \cite[Theorem~2]{fite}. The bound here uses the
adjoint hypothesis and applies to tensor representations as well.
Sheng and Zuo determine the Newton polygons in the Mumford-type
tensor setting, where one active embedding in a nonsplit orbit
can give a binomial distribution of slopes
\cite[Theorem~1.2]{shengzuo}. The formula here is at degree-one
primes over the tensor-defining extension and uses all irreducible
summands of \(H^1\).
